\documentclass[12pt]{article}
\usepackage[margin=1in]{geometry}
\usepackage{amsmath,amssymb}
\usepackage{algorithm,algorithmic}

\title{Problem 93: Arithmetic Expressions}
\author{Project Euler Solution}
\date{}

\begin{document}
\maketitle

\section{Problem Statement}

Choose four distinct digits
\[
\{a,b,c,d\}\subset\{1,2,\dots,9\},\qquad a<b<c<d.
\]
Using each digit exactly once, together with the operations $+,-,\times,\div$ and arbitrary parentheses, we obtain a set of values. We seek the digit set for which the obtainable positive integers contain the longest consecutive run
\[
1,2,\dots,n.
\]
The required output is the concatenation $abcd$.

\section{Exact rational search}

Every valid expression evaluates to a rational number, so the search can be carried out exactly in $\mathbb{Q}$.

For a multiset $S$ of rational numbers, let $R(S)$ be the set of all rational values obtainable by combining the elements of $S$ with the four binary operations, using every element of $S$ exactly once.

\subsection*{Recursive characterization}

If $|S|=1$, then $R(S)=S$.

If $|S|\ge 2$, choose two distinct elements $x,y\in S$, remove them, and replace them by one of
\[
x+y,\qquad x-y,\qquad y-x,\qquad xy,\qquad x/y\ (y\ne 0),\qquad y/x\ (x\ne 0).
\]
Then recurse on the smaller multiset.

This recursion is complete. Any recursive contraction sequence clearly corresponds to a fully parenthesized expression. Conversely, every full binary expression tree has a lowest internal node whose two children are leaves; evaluating that node first contracts two operands into one value and reduces the number of leaves by one. Repeating this process yields a recursive contraction sequence. Thus the recursion generates exactly all obtainable values.

\section{Complexity}

For a multiset of size $m$, there are at most $\binom{m}{2}$ choices of unordered pair, and for each pair at most $6$ legal replacement values. Hence the number of terminal branches for one digit set is bounded by
\[
6\binom{4}{2}\cdot 6\binom{3}{2}\cdot 6\binom{2}{2}
=3888.
\]

There are
\[
\binom{9}{4}=126
\]
possible digit sets, so the total search is very small.

\section{Algorithm}

\begin{algorithm}[H]
\caption{Search all 4-digit sets}
\begin{algorithmic}[1]
\STATE $\mathrm{best\_run}\gets 0$, $\mathrm{best\_digits}\gets \text{``''}$
\FOR{each 4-element subset $D\subset\{1,\dots,9\}$}
    \STATE Compute all exact rational values reachable from $D$ by recursive pair contraction
    \STATE Keep only positive integers
    \STATE Compute the largest $n$ such that $1,2,\dots,n$ are all reachable
    \IF{$n>\mathrm{best\_run}$}
        \STATE Update $\mathrm{best\_run}$ and $\mathrm{best\_digits}$
    \ENDIF
\ENDFOR
\RETURN $\mathrm{best\_digits}$
\end{algorithmic}
\end{algorithm}

\section{Result}

The optimal set is
\[
\{1,2,5,8\},
\]
which yields every positive integer from $1$ through $51$. Therefore the required concatenation is
\[
\boxed{1258}.
\]

\section{Answer}

\[
\boxed{1258}
\]

\end{document}
